\chapter{Chapter06 — Viola Seed Length}
%%% generated by the blueprint pipeline from FABL.Chapter06.FoolingF₂Polynomials.ViolaSeedLength
\section{Overview}

Declarations of the Lean module \texttt{FABL.Chapter06.FoolingF₂Polynomials.ViolaSeedLength}.

\section{Declarations}

\begin{definition}[Base bias parameter]
\lean{FABL.violaBaseBias}
\statementsource{Analysis-of-Boolean-Functions-by-Ryan-ODonnell}{pdf-ffbe9c7b3306-p144-b007}
\label{FABL.violaBaseBias}
\leanok
For a real number $\varepsilon$ and a natural number $d$, define the base bias
$\beta_d(\varepsilon)$ by $\beta_d(\varepsilon)=(\varepsilon/9)^{2^{d-1}}$, where the
exponent is interpreted using natural subtraction in $d-1$.
\end{definition}

\begin{theorem}[Positivity of the base bias]
\lean{FABL.violaBaseBias_pos}
\label{FABL.violaBaseBias_pos}
\leanok
\uses{FABL.violaBaseBias}
\difficulty{3}
For every real number $\varepsilon>0$ and every natural number $d$, the quantity
$(\varepsilon/9)^{2^{\max(d-1,0)}}$ is positive.
\end{theorem}

\begin{theorem}[Upper bound for the base bias]
\lean{FABL.violaBaseBias_le_half}
\label{FABL.violaBaseBias_le_half}
\leanok
\uses{FABL.violaBaseBias}
\difficulty{4}
For every real number $\varepsilon$ with $0\le\varepsilon\le1$ and every natural number
$d$, the quantity $(\varepsilon/9)^{2^{\max(d-1,0)}}$ is at most $1/2$.
\end{theorem}

\begin{theorem}[Error at the base bias]
\lean{FABL.violaError_baseBias}
\label{FABL.violaError_baseBias}
\leanok
\uses{FABL.violaBaseBias, FABL.violaError}
\difficulty{4}
For every real number $\varepsilon>0$ and natural number $d\ge1$, set
$b=(\varepsilon/9)^{2^{d-1}}$. Then $9b^{1/2^{d-1}}=\varepsilon$.
\end{theorem}

\begin{definition}[Small bias seed length]
\lean{FABL.violaSmallBiasSeedBits}
\statementsource{Analysis-of-Boolean-Functions-by-Ryan-ODonnell}{
  pdf-ffbe9c7b3306-p153-b002,
  pdf-ffbe9c7b3306-p153-b003,
  pdf-ffbe9c7b3306-p153-b004
}
\label{FABL.violaSmallBiasSeedBits}
\leanok
For natural numbers $d$ and $\ell$, define the seed length to be $2d\ell$ bits.
\end{definition}

\begin{definition}[Small bias sampler]
\lean{FABL.violaSmallBiasSampler}
\statementsource{Analysis-of-Boolean-Functions-by-Ryan-ODonnell}{
  pdf-ffbe9c7b3306-p153-b002,
  pdf-ffbe9c7b3306-p153-b003,
  pdf-ffbe9c7b3306-p153-b004
}
\label{FABL.violaSmallBiasSampler}
\leanok
\uses{FABL.BinaryExtensionField, FABL.F₂Cube, FABL.smallBiasGenerator}
Let $n,d,\ell$ be natural numbers with $\ell\ne 0$, and choose for each of the $d$
indices a pair of elements of the binary extension field of degree $\ell$ over $\mathbb
F_2$. The sampler outputs the sum in $\mathbb F_2^n$ of the $d$ vectors produced by the
small-bias generator from the corresponding field-element pairs.
\end{definition}

\begin{theorem}[Size of the sampler seed space]
\lean{FABL.card_violaSmallBiasSamplerSeedSpace}
\label{FABL.card_violaSmallBiasSamplerSeedSpace}
\leanok
\uses{FABL.BinaryExtensionField, FABL.binaryExtensionField_natCard, FABL.violaSmallBiasSeedBits}
\difficulty{3}
Let $d$ be a natural number and let $\ell$ be a nonzero natural number. The set of
$d$-tuples of pairs of elements of the degree-$\ell$ extension field of the two-element
field has cardinality $2^{2d\ell}$.
\end{theorem}

\begin{definition}[Viola seed logarithmic bound]
\lean{FABL.violaSeedLogBound}
\statementsource{Analysis-of-Boolean-Functions-by-Ryan-ODonnell}{pdf-ffbe9c7b3306-p153-b002}
\label{FABL.violaSeedLogBound}
\leanok
For natural numbers $n$ and $d$ and a real number $\varepsilon$, define the Viola seed
logarithmic bound to be $d\bigl(4+2(\log_2 n+2^{d-1}\log_2(9/\varepsilon))\bigr)$, where
$d-1$ in the exponent is natural-number subtraction.
\end{definition}

\begin{theorem}[Logarithmic bias support identity]
\lean{FABL.logb_smallBiasSupportEnvelope_eq}
\label{FABL.logb_smallBiasSupportEnvelope_eq}
\leanok
\uses{FABL.violaBaseBias, FABL.violaBaseBias_pos}
\difficulty{4}
Let $n,d$ be natural numbers and let $\varepsilon>0$, with $n>0$. Define
$b=(\varepsilon/9)^{2^{d-1}}$, where $d-1$ is truncated natural subtraction. Then
\[\log_2\!\left(16\left(\frac{n}{b}\right)^2\right)=4+2\left(\log_2 n+2^{d-1}\log_2\!\left(\frac{9}{\varepsilon}\right)\right).\]
\end{theorem}

\begin{theorem}[Seed length from generator size]
\lean{FABL.violaSmallBiasSeedBits_le_of_card}
\label{FABL.violaSmallBiasSeedBits_le_of_card}
\leanok
\uses{FABL.logb_smallBiasSupportEnvelope_eq, FABL.smallBiasGeneratorMultiset, FABL.smallBiasGeneratorMultiset_card, FABL.violaBaseBias, FABL.violaSeedLogBound, FABL.violaSmallBiasSeedBits}
\difficulty{5}
Let $n,d,\ell$ be natural numbers with $n>0$ and $\ell\neq0$, and let $\varepsilon>0$.
For the small-bias generator formed from pairs in the degree-$\ell$ extension field of
$\mathbb F_2$, suppose the cardinality of its output multiset is at most
$16\bigl(n/B(\varepsilon,d)\bigr)^2$, where
$B(\varepsilon,d)=(\varepsilon/9)^{2^{d-1}}$. Then its seed length $2d\ell$ is at most
$d\bigl(4+2(\log_2 n+2^{d-1}\log_2(9/\varepsilon))\bigr)$.
\end{theorem}

\begin{theorem}[Exists viola small bias distribution]
\lean{FABL.exists_violaSmallBiasDistribution}
\statementsource{Analysis-of-Boolean-Functions-by-Ryan-ODonnell}{
  pdf-ffbe9c7b3306-p153-b002,
  pdf-ffbe9c7b3306-p153-b003,
  pdf-ffbe9c7b3306-p153-b004
}
\label{FABL.exists_violaSmallBiasDistribution}
\leanok
\uses{FABL.ProbabilityDensity.Fools, FABL.ProbabilityDensity.IsBiased.violaTheorem, FABL.ProbabilityDensity.convolutionPower, FABL.exists_smallBiasGenerator_of_real, FABL.f₂PolynomialSignClass, FABL.smallBiasGeneratorDensity, FABL.violaBaseBias, FABL.violaBaseBias_le_half, FABL.violaBaseBias_pos, FABL.violaError_baseBias, FABL.violaSeedLogBound, FABL.violaSmallBiasSeedBits, FABL.violaSmallBiasSeedBits_le_of_card}
\difficulty{4}
% dataset-pipeline: stub -- prose pending build_informal_statements.py
Statement pending regeneration.
\end{theorem}
